\section{Introducción}\label{sec:intro}

El objetivo del proyecto es predecir si un cliente tendrá dificultades de pago el mes siguiente, variable \var{payment_difficulty_next_month}. La Fase I se limitó deliberadamente al modelo lógico, un árbol de decisión, con el fin de obtener reglas que permitieran entender qué perfiles eran asignados a cada clase. La Fase II amplía ese propósito y compara tres familias de modelos (árbol, regresión logística y modelo aditivo explicable, EBM) no solo por su capacidad de discriminación, sino por el tipo de explicaciones globales y locales que cada una hace posibles. Por tanto, el objetivo no es únicamente obtener una predicción, sino justificar el compromiso entre rendimiento y legibilidad.

El análisis se realiza sobre el conjunto de entrenamiento, que contiene 25\,000 observaciones y 15 predictores, 11 numéricos y 4 categóricos. El conjunto de test, con 5\,000 filas y sin etiquetas, se reserva para la predicción final y no interviene en el desarrollo. Los datos provienen de \emph{Default of Credit Card Clients} \cite{uci}. Dado que la clase positiva es minoritaria, con 5\,303 casos (21,21\,\%) frente a 19\,697 (78,79\,\%), la evaluación debe considerar el comportamiento en ambas clases. Predecir siempre la clase mayoritaria bastaría para alcanzar un 78,79\,\% de exactitud, una cifra que por sí sola no informa sobre la detección de dificultades de pago.

Para que la comparación sea homogénea, las tres familias se evalúan sobre las mismas cinco particiones externas, agrupadas por perfil idéntico de los quince predictores, con validación cruzada anidada dentro de cada entrenamiento. La agrupación evita que perfiles repetidos, 31 en total y 13 con etiquetas distintas, aparezcan simultáneamente en entrenamiento y validación; el anidamiento separa la selección de hiperparámetros de la estimación del rendimiento. La configuración se elige por ROC-AUC medio y el umbral por exactitud equilibrada medio. La comparación final emplea ROC-AUC, exactitud (accuracy) equilibrada, F1, sensibilidad (recall), precisión y exactitud, con la clase positiva igual a uno. El test permanece intacto hasta el ajuste del modelo final.

Los resultados muestran una jerarquía consistente entre particiones. El EBM alcanza la mayor ROC-AUC ($0{,}7755 \pm 0{,}0077$) y la mejor exactitud equilibrada con umbral interno (0,6997), seguido de la regresión logística ($0{,}7515 \pm 0{,}0075$) y del árbol ($0{,}7372 \pm 0{,}0044$). Los umbrales internos mejoran la detección de la clase positiva a costa de precisión y exactitud. Con 0,5 los tres modelos superan la referencia mayoritaria, pero detectan menos del 30\,\% de los positivos. El modelo final es un EBM sin interacciones, con \var{max_bins}=32, un mínimo de 10 observaciones por hoja y umbral 0,20. Aplicado al test, asigna 3\,208 casos a la clase 0 (64,16\,\%) y 1\,792 a la clase 1 (35,84\,\%); sin etiquetas reales no pueden calcularse métricas sobre esa salida.

El informe describe a continuación el método y el protocolo de evaluación, presenta los resultados comparados, analiza las explicaciones globales y locales de cada familia, discute sus implicaciones y límites, y cierra con las conclusiones.
